\clearpage\fancyhead[L]{\small JUNO.}\hypertarget{native-375-lower}{}\label{native-375-lower}
\begin{center}\textbf{JUNO.}\par\vspace{4pt}\hrule\vspace{4pt}\fontsize{10}{12}\selectfont Monatliche Correspondenz. Band XXVIII. Seite 574\;..\;578. Dezember 1813.\par\vspace{4pt}\hrule\vspace{4pt}\end{center}


\begin{flushright}Göttingen, 8., 9. und 16. Dezember 1813.\end{flushright}
\noindent --- --- --- Das Wetter war wie gewöhnlich zu dieser Jahreszeit sehr ungünstig für Beobachtungen. Ich bestimme die Zeit jetzt anhand von Sternazimuten, die ich (am Tag) mit dem Theodolit beobachte. Beobachtungen bei gleichen Sonnenhöhen gelingen jetzt nie, \SourceLeaf{376}da es sich nur gelegentlich für einzelne Stunden etwas aufklart. Gegen Sonnenuntergang vergleiche ich jetzt $\alpha$ Aurigae mit einem 4000 Meter entfernten Turm, dessen Azimut im Sommer bestimmt wurde. Das geht recht gut.

\Needspace{12\baselineskip}\par\indent Ich habe Juno nur ein einziges Mal beobachten können, doch die Beobachtung selbst ist sehr gut und wurde am Tag der Opposition gemacht. Der Planet hatte die achte Größe.
\begin{center}\fontsize{10}{12}\selectfont\setlength{\tabcolsep}{4pt}\renewcommand{\arraystretch}{1.15}\begin{tabular}{l|c|c|c}
1813 & \shortstack{Mittlere Zeit\\in Göttingen} & \shortstack{Rektaszension $\JunoSymbol$} & \shortstack{Südliche\\Deklination $\JunoSymbol$}\\\hline
Nov. 19 & $9^{\mathrm u}\;46^{\mathrm m}\;46^{\mathrm s}$ & $60^\circ\;30'\;35''0$ & $3^\circ\;5'\;9''3$\\
\end{tabular}\end{center}
\Needspace{7\baselineskip}\par\indent Die scheinbare Position des Vergleichssterns wurde aus der \emph{Hist. Cél.} bestimmt
\begin{center}Rektasz. $=60^\circ\;47'\;50''0$\quad Dekl. $=3^\circ\;0'\;16''0$ südl.\end{center}

\Needspace{12\baselineskip}\par\indent Ich habe ohne jedes Bedenken aus dieser Beobachtung die Opposition abgeleitet, da von anderswo wohl kaum Beobachtungen eintreffen werden.
\begin{center}\emph{Achte Opposition der Juno:}\end{center}
\begin{center}1813 November 19. $18^{\mathrm u}\;11^{\mathrm m}\;25^{\mathrm s}$ Mittlere Zeit in Göttingen\end{center}
\begin{center}\fontsize{10}{12}\selectfont\setlength{\tabcolsep}{2pt}\renewcommand{\arraystretch}{1.15}\begin{tabular}{lcr}
Wahre Länge & $.\;.\;.\;.\;.\;.\;.\;.\;.\;.\;.\;.\;.$ & $57^\circ\;33'\;58''0$\\
Geozentrische Breite & $.\;.\;.\;.\;.\;.\;.\;.\;.\;.$ & $23\;18\;46.8\ \text{südl.}$\\
\end{tabular}\end{center}
\Needspace{6\baselineskip}\begin{center}\emph{Fortsetzung vom 9. Dezember $10^{\mathrm u}$.}\end{center}
\par\indent Heute Abend habe ich bei dichtem Nebel und Mondschein zwar noch nicht wieder an der Stelle nachgesehen, an der ich Juno am 19. November beobachtet habe, aber ich habe Juno selbst wieder aufgesucht. \emph{Verhältnismäßig} schien sie heute \emph{heller} als am 19. November. Wenigstens war sie deutlich heller als der Stern achter Größe, mit dem sie heute verglichen wurde, während sie am 19. November den damals zum Vergleich herangezogenen Stern achter Größe kaum erreichte. Ich bitte Herrn \textsc{Nicolai}, die Beobachtungen zu reduzieren und mit den unten folgenden Elementen zu vergleichen. Die Beobachtung beruht zwar nur auf drei kümmerlichen Vergleichen mit einem Stern aus der \emph{Hist. Cél.}, der am 20. Januar 1798 am zweiten Faden um $3^{\mathrm u}\;42^{\mathrm m}\;48^{\mathrm s}5$ beobachtet wurde. Sie kann aber dennoch zur Kontrolle dienen.
\begin{center}
1813 Dezember 9. $0^{\mathrm u}\;22^{\mathrm m}\;2^{\mathrm s}$ Sternzeit, $\JunoSymbol$ folgt auf den Stern
$3^{\mathrm m}\;49^{\mathrm s}75$ Zeit $=57'\;26''25$ Bogen

$0^{\mathrm u}\;32^{\mathrm m}\;4^{\mathrm s}$ --- $\JunoSymbol$ steht weiter südlich um $5'\;0''9$.
\end{center}
\noindent Ungefähr folgt daraus
\begin{center}Rektasz. $=56^\circ\;56\tfrac12'$\quad Dekl. $=3^\circ\;56\tfrac12'$ südl.\end{center}

\Needspace{15\baselineskip}\par\indent Die Elemente, die auf den Oppositionen von 1810, 1811, 1812, 1813 beruhen, lauten (X):
\begin{center}\fontsize{9}{11}\selectfont\setlength{\tabcolsep}{2pt}\renewcommand{\arraystretch}{1.15}\begin{tabular}{lcr}
Epoche 1814, Meridian von Göttingen & $.\;.\;.\;.\;.\;.$ & $65^\circ\;19'\;23''8$\\
Tägliche tropische Bewegung & $.\;.\;.\;.\;.\;.\;.\;.\;.$ & $812''709$\\
Log. der großen Halbachse & $.\;.\;.\;.\;.\;.\;.\;.\;.\;.$ & $0.4267631$\\
Perihel (1814) & $.\;.\;.\;.\;.\;.\;.\;.\;.\;.\;.\;.\;.\;.$ & $53^\circ\;10'\;0''40$\\
Exzentrizitätswinkel & $.\;.\;.\;.\;.\;.\;.\;.\;.\;.\;.\;.\;.$ & $14\;43\;7.97$\\
Aufsteigender Knoten (1814) & $.\;.\;.\;.\;.\;.\;.\;.\;.\;.$ & $171\;10\;33.00$\\
Neigung & $.\;.\;.\;.\;.\;.\;.\;.\;.\;.\;.\;.\;.\;.\;.\;.\;.$ & $13\;4\;12.03$\\
\end{tabular}\end{center}
\par\indent Der Vergleich dieser Elemente, besonders der mittleren Bewegung, mit den Elementen (VIII), die auf den Oppositionen von 1804, 1806, 1807, 1808 beruhten, zeigt, wie stark sich die Jupiterstörungen auswirken.

\Needspace{18\baselineskip}\begin{flushright}\leavevmode\SourceLeaf{377}Göttingen, 16. Dezember 1813.\end{flushright}
\noindent --- --- --- Ich habe mir doch die Mühe gemacht, die Elemente auch an die letzte Opposition der Pallas, die IX., anzupassen. Die Korrekturen sind äußerst gering:
\begin{center}\fontsize{10}{12}\selectfont\setlength{\tabcolsep}{2pt}\renewcommand{\arraystretch}{1.15}\begin{tabular}{lcr}
Länge des Knotens & $.\;.\;.\;.\;.\;.$ & $-2''7$\\
Neigung der Bahn & $.\;.\;.\;.\;.\;.$ & $-1.4$\\
Länge des Perihels & $.\;.\;.\;.\;.\;.$ & $-9.4$\\
Exzentrizitätswinkel & $.\;.\;.\;.\;.\;.$ & $+1.8$\\
Epoche 1810 & $.\;.\;.\;.\;.\;.\;.$ & $-4.2$\\
Tägliche Bewegung & $.\;.\;.\;.$ & $-0''00569$\\
Log. der großen Halbachse & $.$ & $+0.0000021$\\
\end{tabular}\end{center}
\Needspace{22\baselineskip}\par\indent Die 9 Oppositionen stimmen nun folgendermaßen mit den verbesserten Elementen überein:
\begin{center}Unterschied:\end{center}
\begin{center}\fontsize{10}{12}\selectfont\setlength{\tabcolsep}{7pt}\renewcommand{\arraystretch}{1.15}\begin{tabular}{l|c|c}
Oppos. & \shortstack{Mittlere\\Länge} & \shortstack{Heliozentr.\\Breite.}\\\hline
I & $-1''3$ & $-5''3$\\
II & $-12.9$ & $-11.9$\\
III & $-20.4$ & $-0.4$\\
IV & $+2.5$ & $-6.8$\\
V & $+23.5$ & $+9.0$\\
VI & $+21.5$ & $-3.1$\\
VII & $+14.1$ & $-7.2$\\
VIII & $-16.1$ & $+4.0$\\
IX & $-10.9$ & $-0.0$\\
X & $-14.0$ & $+2.4$\\
\end{tabular}\end{center}
\par\indent Ich bin sehr neugierig, wie die Unterschiede ausfallen werden, sobald die Saturnstörungen hinzukommen. [Die Angabe für Oppos. X ist eine handschriftliche Bemerkung. \textsc{Schering}.]
